% !TEX root = ../main.tex

\documentclass[../main.tex]{subfiles}

\graphicspath{{\subfix{../images/}}}

\begin{document}

\subsection{Обзор предметной области}

Современный человек сталкивается с большим количеством личных и профессиональных целей. Для поддержания мотивации и систематизации усилий всё шире применяются методы геймификации~--- использование игровых механик (очки, уровни, награды, прогресс-бары) в неигровых контекстах. Геймификация помогает превращать абстрактные намерения в конкретные, измеримые шаги и повышает вовлечённость пользователя в долгосрочных задачах.

Отдельное направление составляют трекеры привычек и достижений~--- мобильные и веб-приложения, позволяющие фиксировать выполнение повторяющихся действий или достижение заранее определённых целей. Такие решения варьируются от простых чек-листов до сложных RPG-систем с социальным взаимодействием. Несмотря на разнообразие предложений на рынке, пользователи часто сталкиваются с рядом типичных проблем:
\begin{itemize}[leftmargin=*]
	\item отсутствие структурированного представления целей в виде коллекций (<<паков>>) с пошаговым прогрессом;
	\item невозможность делиться готовыми наборами достижений с другими пользователями;
	\item перегруженный игровой интерфейс либо, напротив, недостаток специализированных функций;
	\item зависимость от одной платформы или нестабильная работа кроссплатформенных решений.
\end{itemize}

Мобильные устройства и современные кроссплатформенные технологии разработки позволяют создавать специализированные приложения, объединяющие учёт достижений, визуализацию прогресса и синхронизацию данных между устройствами. Разрабатываемое приложение \textbf{Achi} направлено на решение указанных задач: пользователь может собирать коллекции достижений, отслеживать выполнение шагов, создавать собственные <<паки>> и добавлять в коллекцию наборы, созданные другими участниками.

\subsection{Анализ существующих решений}

В данном разделе представлен анализ современных решений для мониторинга привычек, задач и достижений. Рассмотрены как специализированные приложения, так и универсальные инструменты продуктивности. Особое внимание уделяется функциональным возможностям, качеству пользовательского интерфейса и пригодности для сценария <<структурированные коллекции достижений с пошаговым прогрессом>>.

\subsubsection{Habitica}

Habitica~--- популярное приложение для формирования привычек и отслеживания задач, использующее элементы ролевых игр (RPG)~\cite{habitica}. Пользователь выполняет привычки, ежедневные задачи и пункты списка дел, получая за это опыт, золото и экипировку для игрового персонажа. Предусмотрено объединение в группы для совместного выполнения квестов.

\textbf{Преимущества:} высокий уровень вовлечённости благодаря геймификации; развитое сообщество; бесплатный базовый функционал.

\textbf{Недостатки:} перегруженный игровой интерфейс, отвлекающий от реальных задач; отсутствие концепции <<паков>> достижений с пошаговым прогрессом; сложность адаптации для серьёзных или профессиональных целей.

\subsubsection{Strides}

Strides~--- трекер привычек и целей с акцентом на аналитику и визуализацию прогресса~\cite{strides}. Поддерживаются четыре типа целей: привычка, цель, среднее значение и проект. Приложение предоставляет подробную статистику, графики и напоминания.

\textbf{Преимущества:} качественная визуализация данных; гибкая настройка целей; стабильный интерфейс.

\textbf{Недостатки:} ограниченный бесплатный функционал (лимит на количество целей); отсутствие социальной составляющей и возможности делиться наборами целей; преимущественная ориентация на платформу iOS.

\subsubsection{Any.do}

Any.do~--- менеджер задач и календарь с поддержкой списков, напоминаний и голосового ввода~\cite{anydo}. Приложение ориентировано на повседневное планирование и управление делами.

\textbf{Преимущества:} простой и понятный интерфейс; кроссплатформенность; интеграция с календарём.

\textbf{Недостатки:} отсутствие специализации на достижениях; нет концепции коллекций или пошагового прогресса внутри <<награды>>; ограниченные возможности совместного использования пользовательских наборов задач.

\subsubsection{Achievement Box}

Achievement Box~--- приложение с геймифицированным подходом к формированию полезных привычек~\cite{achievement-box}. За выполнение <<хороших>> привычек пользователь зарабатывает монеты, которые можно тратить в магазине наград. Предусмотрена статистика, графики и резервное копирование через сервисы Google.

\textbf{Преимущества:} стабильная работа и корректное отображение интерфейса; открытый исходный код; развитая система личной статистики.

\textbf{Недостатки:} достижения смешаны с понятием <<привычек>>; отсутствует возможность передавать созданные наборы другим пользователям; нет профиля и социальных функций; ограниченный выбор иконок (нельзя загрузить собственные изображения).

\subsubsection{The Achievements}

The Achievements~--- приложение для создания, отправки и получения достижений с социальными функциями~\cite{the-achievements}. Поддерживается генерация достижений с помощью ИИ, отправка через QR-коды и ссылки, а также лайки, комментарии и подписки.

\textbf{Преимущества:} развитая система пользователей и профилей; возможность выдавать достижения другим; кастомизация стилей.

\textbf{Недостатки:} проблемы с интерфейсом (некорректные отступы, артефакты вёрстки); полная зависимость от интернета; нельзя <<взять>> достижение для самостоятельного выполнения~--- решение принимает только создатель; отсутствует пошаговый прогресс внутри достижения.

\subsubsection{Notion и Todoist}

Notion~--- универсальный инструмент для управления задачами, заметками и проектами~\cite{notion}. Пользователь может организовать информацию в виде страниц, списков и таблиц, настроив собственную систему учёта прогресса. Todoist~--- менеджер задач с поддержкой проектов, меток и фильтров~\cite{todoist}, ориентированный на учёт дедлайнов и повседневных дел.

\textbf{Преимущества Notion:} высокая степень персонализации; подходит для сложных проектов; совместная работа.

\textbf{Недостатки Notion:} сложность освоения; отсутствие встроенной геймификации и специализированного интерфейса <<достижений>>; требуется ручная настройка всех таблиц и связей.

\textbf{Преимущества Todoist:} удобное управление задачами; интеграция с календарями; кроссплатформенность.

\textbf{Недостатки Todoist:} не специализируется на <<достижениях>> или <<коллекциях>>; отсутствует визуальное представление прогресса в виде наград или бейджей с пошаговой структурой.

Результаты сравнения аналогов представлены в табл.~\ref{tab:analogs-compare}.

\aboverulesep=0ex
\belowrulesep=0ex
\begin{small}
\begin{longtable}[t]{
	|>{\RaggedRight\arraybackslash}p{0.18\linewidth}|
	>{\centering\arraybackslash}p{0.09\linewidth}|
	>{\centering\arraybackslash}p{0.09\linewidth}|
	>{\centering\arraybackslash}p{0.09\linewidth}|
	>{\centering\arraybackslash}p{0.09\linewidth}|
	>{\centering\arraybackslash}p{0.09\linewidth}|
	>{\centering\arraybackslash}p{0.09\linewidth}|
	>{\centering\arraybackslash}p{0.09\linewidth}|
	>{\centering\arraybackslash}p{0.09\linewidth}|
}

\caption{Сравнительный анализ аналогов} \label{tab:analogs-compare}\\
\hline
Критерий & Habitica & Strides & Any.do & Achiev. Box & The Achiev. & Notion & Todoist & Achi \\
\hline
\endfirsthead

\caption*{Продолжение табл. \ref{tab:analogs-compare}\raggedleft}\\
\hline
Критерий & Habitica & Strides & Any.do & Achiev. Box & The Achiev. & Notion & Todoist & Achi \\
\hline
\endhead

\endfoot
\hline
\endlastfoot

Создание своих достижений & Нет & Нет & Нет & Частично & Да & Вручную & Нет & Да \\
\hline
Шеринг наборов (паков) & Нет & Нет & Нет & Нет & Да & Шаблоны & Нет & Да \\
\hline
Пошаговый прогресс & Нет & Частично & Нет & Нет & Нет & Вручную & Нет & Да \\
\hline
Профиль и социальность & Да & Нет & Нет & Нет & Да & Да & Да & Да \\
\hline
Стабильность UI/UX & Средняя & Высокая & Высокая & Высокая & Низкая & Средняя & Высокая & --- \\
\hline
Кросс\-плат\-фор\-мен\-ность & Да & Частично & Да & Android & Да & Да & Да & Да \\
\hline
Геймификация & Да & Нет & Нет & Да & Нет & Нет & Частично & Частично \\

\end{longtable}
\end{small}

Таким образом, существующие аналоги обладают разными преимуществами и недостатками. Большинство из них либо не предоставляют возможности создавать структурированные коллекции (<<паки>>) достижений с пошаговым прогрессом, либо не поддерживают их передачу другим пользователям. Популярные трекеры привычек (Habitica, Strides) и таск-менеджеры (Todoist, Any.do) не сфокусированы на концепции <<достижений>> как самостоятельной сущности. Это подтверждает актуальность разработки специализированного кроссплатформенного приложения Achi.

\subsection{Формулировка требований к приложению}

На основе анализа предметной области и существующих решений сформулированы требования к разрабатываемому приложению Achi.

\subsubsection{Пользовательские функции}

\textbf{Управление учётной записью:}
\begin{itemize}[leftmargin=*]
	\item регистрация и авторизация пользователя для синхронизации прогресса и коллекций паков с сервером;
	\item сохранение сессии (токен авторизации) и возможность выхода из аккаунта.
\end{itemize}

\textbf{Работа с достижениями и паками:}
\begin{itemize}[leftmargin=*]
	\item просмотр списка паков в коллекции пользователя;
	\item детализация достижений внутри пака: описание, изображения, текущий прогресс;
	\item отслеживание прогресса: отметка выполнения отдельных шагов (подзадач) с автоматическим расчётом общего прогресса и синхронизацией с сервером;
	\item добавление паков по уникальному коду, созданных другими пользователями;
	\item создание собственных паков: добавление достижений, настройка шагов, загрузка изображений-превью;
	\item редактирование, копирование и удаление созданных паков.
\end{itemize}

\textbf{Интерфейс и навигация:}
\begin{itemize}[leftmargin=*]
	\item нижняя панель навигации (Bottom Nav) с разделами <<Добавление>>, <<Мои достижения>>, <<Профиль>>;
	\item локализация интерфейса (русский и английский языки);
	\item просмотр пользовательского соглашения и политики конфиденциальности.
\end{itemize}

\subsubsection{Системные функции}

\begin{itemize}[leftmargin=*]
	\item синхронизация данных: фоновая отправка прогресса шагов на сервер с оптимистичным обновлением пользовательского интерфейса;
	\item взаимодействие с REST API через сгенерированные клиенты (OpenAPI);
	\item отладочная панель для переключения хоста API и быстрой авторизации (только в отладочных сборках).
\end{itemize}

\subsubsection{Нефункциональные требования}

\textbf{Архитектура и производительность:}
\begin{itemize}[leftmargin=*]
	\item построение приложения по паттерну MVVM с выделением слоя Repository;
	\item оптимизированная загрузка и кэширование изображений.
\end{itemize}

\textbf{Кроссплатформенность:}
\begin{itemize}[leftmargin=*]
	\item единая кодовая база для бизнес-логики и пользовательского интерфейса;
	\item поддержка платформ: Android (основная), Web/WasmJS (вторичная), iOS.
\end{itemize}

\textbf{Пользовательский интерфейс:}
\begin{itemize}[leftmargin=*]
	\item использование дизайн-системы Material~3;
	\item адаптивность под различные размеры экранов;
	\item минимальная версия Android~--- API~24 (Android~7.0).
\end{itemize}

\subsection{Выводы по главе}

\begin{enumerate}[leftmargin=*]
	\item Рассмотрена предметная область геймификации и трекинга достижений. Выявлена потребность в специализированном решении, объединяющем структурированные коллекции достижений, пошаговый прогресс и возможность обмена наборами между пользователями.
	\item Проведён сравнительный анализ восьми аналогов (Habitica, Strides, Any.do, Achievement Box, The Achievements, Notion, Todoist). Ни одно из рассмотренных решений не сочетает полностью функции создания <<паков>> с пошаговым прогрессом, шеринга и кроссплатформенной реализации.
	\item На основе анализа сформулированы функциональные и нефункциональные требования к приложению Achi, определяющие ключевые сценарии: регистрация, управление коллекциями, отслеживание прогресса с серверной синхронизацией и создание собственных паков.
\end{enumerate}

\end{document}
